\documentclass[12pt,a4paper]{article}
\usepackage{fontsetup}
\usepackage[english, russian, ukrainian]{babel}

\usepackage{geometry}
\geometry{left=3cm,right=1.5cm,top=1.5cm,bottom=1.5cm}
\usepackage{xcolor}
\usepackage{tabularray}
\UseTblrLibrary{booktabs}
\usepackage[
    colorlinks=true,
    urlcolor = blue,
    linkcolor  = malina,
    citecolor  = green,
    bookmarks = true,
    bookmarksnumbered=true,
    unicode,
    linktoc = all,
    hypertexnames=false,
    pdftoolbar=false,
    pdfpagelayout=TwoPageRight,
    pdfauthor={Ponomarenko S.M. aka sergiokapone},
    pdfdisplaydoctitle=true,
    pdfencoding=auto
    ]{hyperref}

\begin{document}

\begin{center}
    \textbf{\large НАЦІОНАЛЬНИЙ ТЕХНІЧНИЙ УНІВЕРСИТЕТ УКРАЇНИ}\\[0.2em]
    \textbf{\large «КИЇВСЬКИЙ ПОЛІТЕХНІЧНИЙ ІНСТИТУТ\\ імені ІГОРЯ СІКОРСЬКОГО»}
\end{center}

\vspace{1.5em}

\begin{center}
    \textbf{\Large витяг}\\[0.3em]
    \textbf{з протоколу № 13/2026 засідання Вченої ради}\\[0.2em]
    \textbf{Навчально-наукового}\\[0.2em]
    \textbf{Фізико-технічного інституту}\\[0.3em]
    \textbf{від 26 жлвтня 2026 р.}
\end{center}

\vspace{1em}

\noindent На засіданні присутні 15 із 23 членів Вченої ради.

\vspace{1.5em}

\noindent \textbf{СЛУХАЛИ:} Про рекомендацію рукопису українською мовою «Електрика та магнетизм», автора С. М. Пономаренка, для здобувачів ступеня магістр за спеціальністю G8 Матеріалознавство до видання та визнання як підручника.

\vspace{1em}

\noindent \textbf{УХВАЛИЛИ:}
\begin{enumerate}
    \item Рекомендувати рукопис «Рентгеноструктурний аналіз дисперсних матеріалів», автори О. І. Юркова, П. І. Лобода, О. П. Карасевська, І. Ю. Троснікова, до видання.
    \item Авторам подати рукопис до видавництва «Політехніка» для отримання редакційного висновку.
    \item Рукопис та пакет супровідних документів направити на розгляд Експертній раді з навчальних видань для визнання рукопису як підручника.
\end{enumerate}

\vspace{3em}

\noindent
\begin{minipage}[t]{0.65\textwidth}
    Голова вченої ради\\
    НН ІМЗ ім. О. Патона\\
    д. т. н., професор
\end{minipage}%
\begin{minipage}[t]{0.35\textwidth}
    \raggedleft
    Анатолій МІНІЦЬКИЙ
\end{minipage}

\vspace{2.5em}

\noindent
\begin{minipage}[t]{0.65\textwidth}
    Учений секретар Вченої ради\\
    НН ІМЗ ім. Є. О. Патона\\
    доктор філософії
\end{minipage}%
\begin{minipage}[t]{0.35\textwidth}
    \raggedleft
    Михайло ВТЕРКОВСЬКИЙ
\end{minipage}

\end{document}